\chapter{Glossary}
\label{app:glossary}

\begin{description}[leftmargin=!,labelwidth=3.6cm,style=nextline]

\item[Activity cell] The cell where a user spends a long continuous period
  during the working part of the day, as defined by the operator's own
  classification rules. Paired with the home cell, it is what a
  \emph{transition window} is supposed to sit between.

\item[Baseline (persistence)] The rule ``predict that the next \cid{cell\_id} is
  the last one seen''. No model, no parameters, no training. Its value is a
  property of the population, not a constant.

\item[Behavioural group / profile] One of four clusters of users, cut on the rate
  at which their \emph{physical site} changes, measured separately over four
  bands of the day. Named G0 (homebody) to G3 (morning mover) in order of
  increasing mobility.

\item[\cid{cell\_id}] The identifier of one radio cell, e.g.\ \cid{BSOCHO1}. It
  decomposes into a technology letter, a site root and a zone/sector index.

\item[ChatML] A convention for storing a training example as a short
  conversation with roles (system, user, assistant). Used here purely as a file
  format.

\item[Cold start] Predicting for a user about whom little history is available.
  Measured here as accuracy at 30\,\% context fraction.

\item[Context fraction] The share of a user's day given to the model before it
  must predict the rest. Part of the measurement, not of the model: moving it
  moves the baseline too.

\item[Copy (levelled corpus)] A scored position whose target is a record
  duplicated from the one before, produced by stretching a sequence to a fixed
  length. Any rule that repeats is correct on it by construction.

\item[Gap] Model accuracy minus persistence-baseline accuracy, in points, both
  measured on the same users and the same scored positions. The only quantity in
  this report that is comparable across datasets.

\item[Home cell] The cell where a user is found both late at night and early in
  the morning, as defined by the operator's own classification rules.

\item[Hour token] A dedicated vocabulary symbol \cid{<H00>}\dots\cid{<H23>}
  placed before each event, giving the hour of day of that event. Always given,
  never predicted.

\item[LoRA] \emph{Low-Rank Adaptation.} A fine-tuning method that freezes the
  original weights and trains small matrices alongside them, reducing the memory
  needed for training by orders of magnitude.

\item[Paired bootstrap] A significance test in which the same test users are
  resampled for both models being compared, so that the difficulty of the draw
  cancels and only the difference between the models remains.

\item[QLoRA] LoRA applied on top of a base model stored in 4-bit precision,
  which reduces the memory the frozen weights occupy.

\item[Record / event] One (\cid{cell\_id}, timestamp) pair. Not a movement: an
  idle handset still produces records, because a base station re-signals to the
  handsets attached to it roughly every four hours on 2G and every two hours on
  3G.

\item[Site root] The part of a \cid{cell\_id} identifying the physical antenna,
  e.g.\ \cid{KVCHE} in \cid{UKVCHE102}. All mobility in this report is measured
  on the site root, because 12.4\,\% of consecutive-event pairs change
  \cid{cell\_id} without changing place.

\item[Teacher forcing] Scoring position $i$ from the \emph{real} history
  $0 \ldots i-1$ rather than from the model's own earlier predictions, so that a
  single mistake cannot cascade.

\item[Top-$k$] The correct cell is among the model's $k$ highest-ranked
  candidates. Top-1, top-3 and top-5 are the same predictions read at three
  depths.

\item[Top-1 within repertoire] Top-1 restricted to targets that are cells present
  in the user's own history, which is the set the model can in principle choose
  from.

\item[Transition window] A range of hours at which a user is more likely to move
  between their home area and their activity area, and where the training loss is
  therefore multiplied. Guessed globally at first (04--06h, 18--20h) and later
  derived per behavioural group from the data.

\item[Zone / sector] The trailing digits of a \cid{cell\_id}, identifying which
  angular sector of the antenna served the handset. A change of sector alone is
  not a movement.

\end{description}
